% Packages and macros shared by every flavour of the paper (article, revtex, ...).
% A flavour entry file (e.g. article.tex) sets \documentclass, then % Packages and macros shared by every flavour of the paper (article, revtex, ...).
% A flavour entry file (e.g. article.tex) sets \documentclass, then % Packages and macros shared by every flavour of the paper (article, revtex, ...).
% A flavour entry file (e.g. article.tex) sets \documentclass, then % Packages and macros shared by every flavour of the paper (article, revtex, ...).
% A flavour entry file (e.g. article.tex) sets \documentclass, then \input{preamble}.

\usepackage[utf8]{inputenc}
\usepackage[T1]{fontenc}
\usepackage{lmodern}
\usepackage{amsmath}
\usepackage{amssymb}
\usepackage{graphicx}
\usepackage{booktabs}
\usepackage{placeins} % \FloatBarrier keeps each insight's figures beside its text
\usepackage[numbers,sort&compress]{natbib}
\usepackage[hidelinks]{hyperref}
\usepackage{cleveref}

% Physics notation used across chapters.
\newcommand{\Msun}{M_\odot}
\newcommand{\rhoc}{\rho_c}
\newcommand{\Mmax}{M_{\max}}
\newcommand{\rh}{r_h}
\newcommand{\Dch}{D}
\newcommand{\RGB}{\mathcal{R}^2_{\mathrm{GB}}} % the Gauss--Bonnet invariant (Section 2.1)
\newcommand{\order}[1]{\mathcal{O}\!\left(#1\right)}

% Draft markers, visible in the PDF and searchable in the source: \todo (todonotes,
% inline, grey) marks a drafting gap; \review (inline, orange) marks a claim the
% supervisors must check; \placeholder marks missing title-block data, where a
% todonotes box cannot sit. Uncomment \listoftodos in the flavour to list them all.
\usepackage{todonotes}
\presetkeys{todonotes}{inline,color=gray!20}{}
\newcommand{\review}[1]{\todo[color=orange!30]{\textbf{Review:} #1}}
\newcommand{\placeholder}[1]{\textbf{[#1]}}
.

\usepackage[utf8]{inputenc}
\usepackage[T1]{fontenc}
\usepackage{lmodern}
\usepackage{amsmath}
\usepackage{amssymb}
\usepackage{graphicx}
\usepackage{booktabs}
\usepackage{placeins} % \FloatBarrier keeps each insight's figures beside its text
\usepackage[numbers,sort&compress]{natbib}
\usepackage[hidelinks]{hyperref}
\usepackage{cleveref}

% Physics notation used across chapters.
\newcommand{\Msun}{M_\odot}
\newcommand{\rhoc}{\rho_c}
\newcommand{\Mmax}{M_{\max}}
\newcommand{\rh}{r_h}
\newcommand{\Dch}{D}
\newcommand{\RGB}{\mathcal{R}^2_{\mathrm{GB}}} % the Gauss--Bonnet invariant (Section 2.1)
\newcommand{\order}[1]{\mathcal{O}\!\left(#1\right)}

% Draft markers, visible in the PDF and searchable in the source: \todo (todonotes,
% inline, grey) marks a drafting gap; \review (inline, orange) marks a claim the
% supervisors must check; \placeholder marks missing title-block data, where a
% todonotes box cannot sit. Uncomment \listoftodos in the flavour to list them all.
\usepackage{todonotes}
\presetkeys{todonotes}{inline,color=gray!20}{}
\newcommand{\review}[1]{\todo[color=orange!30]{\textbf{Review:} #1}}
\newcommand{\placeholder}[1]{\textbf{[#1]}}
.

\usepackage[utf8]{inputenc}
\usepackage[T1]{fontenc}
\usepackage{lmodern}
\usepackage{amsmath}
\usepackage{amssymb}
\usepackage{graphicx}
\usepackage{booktabs}
\usepackage{placeins} % \FloatBarrier keeps each insight's figures beside its text
\usepackage[numbers,sort&compress]{natbib}
\usepackage[hidelinks]{hyperref}
\usepackage{cleveref}

% Physics notation used across chapters.
\newcommand{\Msun}{M_\odot}
\newcommand{\rhoc}{\rho_c}
\newcommand{\Mmax}{M_{\max}}
\newcommand{\rh}{r_h}
\newcommand{\Dch}{D}
\newcommand{\RGB}{\mathcal{R}^2_{\mathrm{GB}}} % the Gauss--Bonnet invariant (Section 2.1)
\newcommand{\order}[1]{\mathcal{O}\!\left(#1\right)}

% Draft markers, visible in the PDF and searchable in the source: \todo (todonotes,
% inline, grey) marks a drafting gap; \review (inline, orange) marks a claim the
% supervisors must check; \placeholder marks missing title-block data, where a
% todonotes box cannot sit. Uncomment \listoftodos in the flavour to list them all.
\usepackage{todonotes}
\presetkeys{todonotes}{inline,color=gray!20}{}
\newcommand{\review}[1]{\todo[color=orange!30]{\textbf{Review:} #1}}
\newcommand{\placeholder}[1]{\textbf{[#1]}}
.

\usepackage[utf8]{inputenc}
\usepackage[T1]{fontenc}
\usepackage{lmodern}
\usepackage{amsmath}
\usepackage{amssymb}
\usepackage{graphicx}
\usepackage{booktabs}
\usepackage{placeins} % \FloatBarrier keeps each insight's figures beside its text
\usepackage[numbers,sort&compress]{natbib}
\usepackage[hidelinks]{hyperref}
\usepackage{cleveref}

% Physics notation used across chapters.
\newcommand{\Msun}{M_\odot}
\newcommand{\rhoc}{\rho_c}
\newcommand{\Mmax}{M_{\max}}
\newcommand{\rh}{r_h}
\newcommand{\Dch}{D}
\newcommand{\RGB}{\mathcal{R}^2_{\mathrm{GB}}} % the Gauss--Bonnet invariant (Section 2.1)
\newcommand{\order}[1]{\mathcal{O}\!\left(#1\right)}

% Draft markers, visible in the PDF and searchable in the source: \todo (todonotes,
% inline, grey) marks a drafting gap; \review (inline, orange) marks a claim the
% supervisors must check; \placeholder marks missing title-block data, where a
% todonotes box cannot sit. Uncomment \listoftodos in the flavour to list them all.
\usepackage{todonotes}
\presetkeys{todonotes}{inline,color=gray!20}{}
\newcommand{\review}[1]{\todo[color=orange!30]{\textbf{Review:} #1}}
\newcommand{\placeholder}[1]{\textbf{[#1]}}
